\section{Conclusion}

In this project, our primary goal was to build a machine learning model capable of detecting trees in urban environments. This is crucial for our partner \textit{infrared.city}, as they need precise tree locations to simulate urban temperatures during the summer. We utilized free \textit{Sentinel-2} satellite imagery and open tree register data from various cities to train a specialized \textit{U-Net} model.

The most important result is that our model actually performs very well, even if the formal evaluation metrics initially suggested otherwise. Because the public tree registers only include public trees, our precision score was negatively impacted when the model correctly identified trees in private gardens. However, visual inspections confirmed that the model's predictions are highly accurate. It successfully learned the visual characteristics of a tree from the satellite data and can find them everywhere, not just on public streets. 

Additionally, incorporating imagery from three different seasons proved to be a very effective strategy. Our data analysis revealed that the NDVI patterns change significantly depending on the city. Thanks to the temporal attention layers, the model was able to learn these patterns and autonomously determine which month is the most informative for each specific location. 

The dual-head U-Net architecture was the right choice for this specific problem. Because our urban maps contain many pixels with zero trees, a standard model would often default to predicting zero everywhere to minimize the loss. By splitting the model into two parts---one predicting the probability of a tree's presence, and the other counting the abundance---we achieved much better results and significantly reduced false predictions on buildings or streets.

Naturally, there is still work to be done in the future. We need to further test the model on cities with different climates, such as Barcelona. Ultimately, however, we successfully built a working, scalable prototype. In the future, the tree maps generated by our model can be integrated directly into temperature models to pinpoint exactly where trees provide shade and help cool down the city.
